% ============================================================
%  Wang Zihao — Curriculum Vitae
%  Engine: XeLaTeX  (compile with:  tectonic main.tex)
%  Also works on Overleaf:  upload main.tex, set compiler to XeLaTeX
%  Font: TeX Gyre Heros (ships with TeX Live / Overleaf)
%
%  Layout follows Owen Morgan's LaTeX CV (github.com/opmorgan/cv):
%    · section headings: bold, with a full-width rule underneath
%    · each entry: date right-aligned in a 15% column, then a short
%      vertical tick (\vline) beside the title, then the content
%    · title bold / institution italic / details in \footnotesize
%    · monochrome (no accent colour)
%  Section taxonomy and content are the author's own.
% ============================================================
\documentclass[11pt]{article}

\usepackage[margin=0.95in]{geometry}
\usepackage{fontspec}
\setmainfont{texgyreheros}[
  Extension=.otf,
  UprightFont=*-regular,
  BoldFont=*-bold,
  ItalicFont=*-italic,
  BoldItalicFont=*-bolditalic
]
\usepackage{titlesec}
\usepackage{enumitem}
\usepackage{needspace}
\usepackage[hidelinks]{hyperref}

% ---- section heading: bold + full-width rule under it (Morgan style) ----
\titleformat{\section}
  {\normalfont\large\bfseries}{}{0em}{}[\titlerule]
\titlespacing*{\section}{0pt}{9pt}{5pt}

% ---- run-in bold label (e.g. "Research Interests:") ----
\newcommand{\subrun}[1]{\vspace{2pt}\noindent\textbf{#1}\hspace{0.35em}\ignorespaces}
\newcommand{\cvsec}[1]{\needspace{3\baselineskip}\section*{#1}}

% ---- entry: date | short tick | content ----
% \entry{date}{title}{body}
\newcommand{\entry}[3]{%
  \noindent
  \begin{minipage}[t]{.17\textwidth}\raggedleft #1\end{minipage}%
  \hfill\vline\hfill
  \begin{minipage}[t]{.78\textwidth}
    \textbf{#2}\par
    #3%
  \end{minipage}%
  \par\vspace{.18cm}%
}

% ---- compact single-line entry (same skeleton, tighter) ----
\newcommand{\aline}[2]{%
  \noindent
  \begin{minipage}[t]{.17\textwidth}\raggedleft #1\end{minipage}%
  \hfill\vline\hfill
  \begin{minipage}[t]{.78\textwidth}#2\end{minipage}%
  \par\vspace{.1cm}%
}

% ---- detail bullet list, in \footnotesize (Morgan's "detail note" size) ----
\newenvironment{detail}
  {\footnotesize\begin{itemize}[leftmargin=1.05em,itemsep=0.5pt,topsep=1.5pt,parsep=0pt,label=\textbullet]}
  {\end{itemize}}
\setlist[itemize]{leftmargin=1.05em,itemsep=0.5pt,topsep=1.5pt,parsep=0pt,label=\textbullet}

\pagestyle{empty}
\setlength{\parindent}{0pt}

\begin{document}

% ---------------- header (Morgan style: date left, identity right) -------
\noindent
\begin{minipage}[t]{.38\textwidth}
  \itshape\footnotesize Last updated \today
\end{minipage}%
\hfill
\begin{minipage}[t]{.60\textwidth}
  \raggedleft
  {\huge\bfseries Wang Zihao}

  \vspace{.25em}

  {\small Zhejiang University, Department of Psychology and Behavioral Sciences}

  \vspace{.25em}

  {\footnotesize \href{mailto:psychwangzihao@zju.edu.cn}{psychwangzihao@zju.edu.cn}
   \ \textperiodcentered\ \ +86 139-5729-1709}

  \vspace{.15em}

  {\footnotesize \href{https://github.com/psychwangzihao}{GitHub}
   \ |\ \href{https://orcid.org/0009-0000-3877-3496}{ORCiD}
   \ |\ \href{https://psychwangzihao.github.io}{psychwangzihao.github.io}}
\end{minipage}
\par\vspace{4pt}

% ---------------- education ----------------
\cvsec{Education}
\entry{2025 -- 2029}%
  {B.Sc. in Psychology (Qiushi Honor's Program)}%
  {Zhejiang University, Hangzhou \ \textperiodcentered\ \ expected June 2029
   \par\vspace{1pt}
   \textit{Research mentor: Dr.\ Yuzheng Hu (Department of Psychology and Behavioral Sciences)}
   \par\vspace{2pt}
   \begin{detail}
     \item GPA: 3.95/4.3 (Top 2 in Honor's Program) \ \textperiodcentered\ \ English: IELTS 8.0
     \item Courses taken: Research Techniques in Psychology; Language, Mind and Brain;
           Probability and Mathematical Statistics; Calculus; Linear Algebra; Programming;
           Ethics and Philosophy of Artificial Intelligence
   \end{detail}}

\entry{2026}%
  {Oxford Prospects Winter Visit}%
  {University of Oxford, Oxford \ \textperiodcentered\ \ Jan--Feb 2026
   \par\vspace{1pt}
   \textit{Psychology: Interdisciplinary Frontiers of Behaviour and Mind in the Age of AI (grade A+)}
   \par\vspace{2pt}
   \begin{detail}
     \item Received the Best Presentation Award and the Best Film Award
   \end{detail}}

% ---------------- research interests & training ----------------
\cvsec{Research Interests \& Training}
\subrun{Research Interests:} mental imagery and aphantasia; consciousness (self-model; absence).

\vspace{3pt}
\subrun{Research Training:}
\begin{detail}
  \item \textbf{Lab Rotations:} conducted rotations in social, clinical, developmental and cognitive
        neuroscience labs, progressively focusing on cognitive neuroscience as a pathway to
        investigate consciousness.
  \item \textbf{Methodological \& Philosophical Training:} national academic conferences and
        philosophy-of-mind seminars; Consciousness Science Reading Group (Swarma Club);
        Interdisciplinary Reading Club of Aphantasia (IRCA).
  \item \textbf{Research Participation:} contributed to a study on brain plasticity following
        spinal-cord-stimulation therapy (Liangzhu Laboratory, Zhejiang University), including
        stimulus development and fMRI data acquisition.
  \item \textbf{Technical Skills:} basic training in fMRI, EEG, fNIRS and eye-tracking, with more
        hands-on practice in behavioural experiments and fMRI; Python and C, currently learning
        MATLAB, SPM and PsychoPy.
\end{detail}

% ---------------- research experience ----------------
\cvsec{Research Experience}
\entry{2026 -- 2028}%
  {Undergraduate Researcher, Visual Mental Imagery and Aphantasia}%
  {Zhejiang University, in collaboration with Dr.\ Jianghao Liu (Paris Brain Institute, ICM)
   \par\vspace{2pt}
   \begin{detail}
     \item \textbf{Research question:} whether imagery signals in aphantasia fail to be generated,
           or are generated but fail to reach conscious access --- a distinction the existing
           literature has not resolved.
     \item \textbf{Design \& methods:} co-designed the study; the protocol combines a near-threshold
           perceptual detection task with an imagery task, together with electrophysiological and
           neuromodulation measures of visual-cortical and prefrontal involvement.
     \item \textbf{Role:} built the experimental platform at Zhejiang University; responsible for
           participant recruitment, data collection and quality control.
   \end{detail}}

% ---------------- community service ----------------
\cvsec{Community Service}
\entry{2026 -- Present}%
  {Founder and Coordinator, Consciousness Observers (CO-LAB)}%
  {Zhejiang University
   \par\vspace{2pt}
   \begin{detail}
     \item Founded an interdisciplinary consciousness-research group spanning several departments;
           designed and maintain its website; convened its first event in the Department of
           Psychology.
   \end{detail}}

\entry{2025 -- 2026}%
  {Lead Organizer \& Initiator, Psychology Outreach Series, 2050@2026 Global Youth Gathering}%
  {Hangzhou, China
   \par\vspace{2pt}
   \begin{detail}
     \item Designed and delivered the keynote \emph{Where Is the Mind?}; organized and moderated the
           \emph{This Is Psychology} forum (recruited three peer speakers and co-presented); and
           facilitated the \emph{The Psych Crew Assembles} roundtable.
     \item Coordinated three experiential sessions (nature therapy, movement therapy and an
           interactive exhibition) by managing practitioner invitations, participant logistics and
           on-site support.
   \end{detail}}

\aline{2026 -- 2027}{\textbf{Peer Mentor Group Leader}, Qiushi Psychology Science Class \&
Brain-Machine Intelligence Innovation Class, Zhejiang University}

\aline{2026 -- Present}{\textbf{Class Monitor}, Psychology (Qiushi Science Class) 2501,
Chu Kochen Honors College, Zhejiang University}

\aline{2026}{\textbf{Volunteer}, Sino-European International Conference on Human Cognition
and Artificial Intelligence (Aug--Sep)}

\aline{2026}{\textbf{High-School Science Outreach}, Hangzhou No.\ 2 High School (March);
Yuqian High School, Lin'an (September)}

% ---------------- awards ----------------
\cvsec{Honors \& Awards}
\aline{08/2026}{\textbf{First Prize} (large-project workshop), ``Mind Science and Intelligent
Future'' International Summer School for Top Students in Psychology --- Ministry of Education
Psychology Teaching Steering Committee \& Faculty of Psychology, Beijing Normal University}
\aline{03/2026}{\textbf{Excellent Poster Award}, Psychology (Honor's Program) Research Exhibition,
Zhejiang University}
\aline{04/2025}{\textbf{Best Debater \& Best Team}, Design Intelligence Award (DIA) Debate
Competition}

\vspace{6pt}
\raggedleft
\href{https://github.com/psychwangzihao/cv}{Source code for this CV}

\end{document}
